\section*{Formulas Used in \texttt{count\_expert\_tokens.py}}

This note summarizes the mathematical steps used in
\texttt{MoE\_profiling/count\_expert\_tokens.py}.

\subsection*{1. Setup}

Let
\[
N = \text{num\_of\_experts}, \qquad
G = \text{ep\_size}, \qquad
k = \text{top\_k}, \qquad
B = \text{target batch size}.
\]

The number of experts assigned to each GPU is
\[
n = \frac{N}{G}.
\]

For an expert index \(e \in \{0,1,\dots,N-1\}\), its GPU index is
\[
g(e) = \left\lfloor \frac{e}{n} \right\rfloor.
\]

\subsection*{2. Token-Level Expert Selection}

For one token, the router selects a top-\(k\) expert list
\[
E = (e_1, e_2, \dots, e_k).
\]

This is converted into a GPU routed-load pattern
\[
X = (X_1, X_2, \dots, X_G),
\]
where
\[
X_r = \sum_{j=1}^{k} \mathbf{1}[g(e_j) = r],
\qquad r = 1,2,\dots,G.
\]

By construction,
\[
\sum_{r=1}^{G} X_r = k.
\]

\subsection*{3. Empirical Token Pattern PMF from Source Data}

Suppose the source file contains \(T\) tokens in total, and the GPU pattern
for token \(t\) is \(X^{(t)}\).

The empirical token-level PMF is
\[
\hat{p}(x)
=
\frac{1}{T}
\sum_{t=1}^{T} \mathbf{1}[X^{(t)} = x].
\]

This corresponds to the \texttt{pattern\_pmf} dictionary in the code.

\subsection*{4. Batch-Level GPU Load Vector}

To predict a larger batch of size \(B\), the script assumes that token patterns
are i.i.d. samples from \(\hat{p}(x)\).

Let
\[
X^{[1]}, X^{[2]}, \dots, X^{[B]} \overset{\text{i.i.d.}}{\sim} \hat{p}(x).
\]

Then the total batch-level GPU routed-token load vector is
\[
L = \sum_{b=1}^{B} X^{[b]}
= (L_1, L_2, \dots, L_G).
\]

Its exact distribution is the \(B\)-fold convolution of the token-level PMF:
\[
\mathbb{P}(L = \ell)
=
\sum_{x^{[1]} + \cdots + x^{[B]} = \ell}
\prod_{b=1}^{B} \hat{p}(x^{[b]}).
\]

The code computes this convolution exactly using repeated sparse distribution
convolution with binary exponentiation.

\subsection*{5. Max Routed Tokens per GPU}

The quantity of interest is the maximum routed-token load among GPUs:
\[
Y = \max_{1 \le r \le G} L_r.
\]

The final probability mass function is therefore
\[
\mathbb{P}(Y = m)
=
\sum_{\ell : \max_r \ell_r = m}
\mathbb{P}(L = \ell).
\]

This is the quantity printed in the distribution table:
\[
\mathbb{P}(Y=0), \mathbb{P}(Y=1), \dots, \mathbb{P}(Y=Bk).
\]

\subsection*{6. Expected Value and Percentile}

Once the PMF of \(Y\) is obtained, the expected maximum routed-token load is
\[
\mathbb{E}[Y]
=
\sum_{m=0}^{Bk} m \, \mathbb{P}(Y=m).
\]

The \(q\)-quantile is defined as the smallest \(m\) such that
\[
\sum_{i=0}^{m} \mathbb{P}(Y=i) \ge q.
\]

In particular, the script reports the \(95\)-th percentile:
\[
\mathrm{p95}
=
\min \left\{
m :
\sum_{i=0}^{m} \mathbb{P}(Y=i) \ge 0.95
\right\}.
\]

\subsection*{7. Ideal Average Load per GPU}

The total number of routed assignments in a batch is
\[
Bk.
\]

If they were perfectly balanced across \(G\) GPUs, the ideal average load per GPU would be
\[
\mathrm{avg\_load}
=
\frac{Bk}{G}.
\]

\subsection*{8. Imbalance Ratio}

The script defines the imbalance ratio as
\[
\mathrm{Imbalance\ Ratio}
=
\frac{\mathbb{E}[Y]}{\mathrm{avg\_load}}
=
\frac{\mathbb{E}[Y]}{Bk/G}.
\]

\subsection*{9. Layer-Wise Averaging}

If the model has \(R\) valid layers and each layer yields a PMF
\(\mathbb{P}_{\ell}(Y=m)\), then the final reported average distribution is
\[
\bar{P}(Y=m)
=
\frac{1}{R}
\sum_{\ell=1}^{R} \mathbb{P}_{\ell}(Y=m).
\]

The same averaging is applied to the expected value, percentile, and imbalance ratio:
\[
\overline{\mathbb{E}[Y]}
=
\frac{1}{R}\sum_{\ell=1}^{R}\mathbb{E}_{\ell}[Y].
\]

\[
\overline{\mathrm{p95}}
=
\frac{1}{R}\sum_{\ell=1}^{R}\mathrm{p95}_{\ell}.
\]

\[
\overline{\mathrm{Imbalance\ Ratio}}
=
\frac{1}{R}\sum_{\ell=1}^{R}\mathrm{Imbalance\ Ratio}_{\ell}.
\]
